\documentclass[11pt,a4paper]{article}

% ===== PACKAGES =====
\usepackage[utf8]{inputenc}
\usepackage[T1]{fontenc}
\usepackage[margin=1in]{geometry}
\usepackage{hyperref}
\usepackage{listings}
\usepackage{xcolor}
\usepackage{booktabs}
\usepackage{array}
\usepackage{longtable}
\usepackage{fancyhdr}
\usepackage{titlesec}
\usepackage{tcolorbox}
\usepackage{fontawesome5}
\usepackage{amssymb}
\usepackage{amsmath}
\usepackage{enumitem}
\usepackage{graphicx}

% ===== COLORS =====
\definecolor{codegreen}{rgb}{0.25,0.49,0.48}
\definecolor{codegray}{rgb}{0.5,0.5,0.5}
\definecolor{codepurple}{rgb}{0.58,0,0.82}
\definecolor{codeblue}{rgb}{0.13,0.29,0.53}
\definecolor{codeorange}{rgb}{0.8,0.4,0.0}
\definecolor{backcolour}{rgb}{0.95,0.95,0.95}
\definecolor{warningbg}{rgb}{1.0,0.95,0.85}
\definecolor{warningborder}{rgb}{0.9,0.6,0.2}
\definecolor{tipbg}{rgb}{0.9,0.97,0.9}
\definecolor{tipborder}{rgb}{0.3,0.7,0.3}
\definecolor{notebg}{rgb}{0.9,0.93,1.0}
\definecolor{noteborder}{rgb}{0.3,0.4,0.7}

% ===== C# CODE LISTING STYLE =====
\lstdefinelanguage{CSharp}{
    keywords={abstract,as,base,bool,break,byte,case,catch,char,checked,class,const,continue,decimal,default,delegate,do,double,else,enum,event,explicit,extern,false,finally,fixed,float,for,foreach,goto,if,implicit,in,int,interface,internal,is,lock,long,namespace,new,null,object,operator,out,override,params,private,protected,public,readonly,ref,return,sbyte,sealed,short,sizeof,stackalloc,static,string,struct,switch,this,throw,true,try,typeof,uint,ulong,unchecked,unsafe,ushort,using,var,virtual,void,volatile,while,async,await},
    keywordstyle=\color{codeblue}\bfseries,
    keywords=[2]{Console,AdsClient,AmsNetId,AmsPort,ReadValue,WriteValue,Connect,Disconnect,Exception,Task,CancellationToken},
    keywordstyle=[2]\color{codepurple}\bfseries,
    sensitive=true,
    morecomment=[l]{//},
    morecomment=[s]{/*}{*/},
    commentstyle=\color{codegreen}\itshape,
    stringstyle=\color{codeorange},
    morestring=[b]",
    morestring=[b]',
}

\lstdefinestyle{csharpstyle}{
    language=CSharp,
    backgroundcolor=\color{backcolour},
    basicstyle=\ttfamily\small,
    breakatwhitespace=false,
    breaklines=true,
    captionpos=b,
    keepspaces=true,
    numbers=left,
    numbersep=8pt,
    numberstyle=\tiny\color{codegray},
    showspaces=false,
    showstringspaces=false,
    showtabs=false,
    tabsize=4,
    frame=single,
    framerule=0.5pt,
    rulecolor=\color{codegray},
    xleftmargin=15pt,
    framexleftmargin=15pt,
}

% ===== C++ CODE LISTING STYLE =====
\lstdefinelanguage{CPP}{
    keywords={auto,break,case,char,const,continue,default,do,double,else,enum,extern,float,for,goto,if,int,long,register,return,short,signed,sizeof,static,struct,switch,typedef,union,unsigned,void,volatile,while,class,private,protected,public,virtual,new,delete,this,try,catch,throw,using,namespace,template,typename,bool,true,false,nullptr,constexpr,override,final,noexcept,inline,explicit,mutable,friend},
    keywordstyle=\color{codeblue}\bfseries,
    keywords=[2]{AdsDevice,AdsVariable,AmsNetId,std,string,cout,endl,uint32_t,int32_t,uint16_t},
    keywordstyle=[2]\color{codepurple}\bfseries,
    sensitive=true,
    morecomment=[l]{//},
    morecomment=[s]{/*}{*/},
    commentstyle=\color{codegreen}\itshape,
    stringstyle=\color{codeorange},
    morestring=[b]",
    morestring=[b]',
}

\lstdefinestyle{cppstyle}{
    language=CPP,
    backgroundcolor=\color{backcolour},
    basicstyle=\ttfamily\small,
    breakatwhitespace=false,
    breaklines=true,
    captionpos=b,
    keepspaces=true,
    numbers=left,
    numbersep=8pt,
    numberstyle=\tiny\color{codegray},
    showspaces=false,
    showstringspaces=false,
    showtabs=false,
    tabsize=4,
    frame=single,
    framerule=0.5pt,
    rulecolor=\color{codegray},
    xleftmargin=15pt,
    framexleftmargin=15pt,
}

% ===== TCOLORBOX STYLES =====
\tcbuselibrary{skins,breakable}

\newtcolorbox{warningbox}{
    colback=warningbg,
    colframe=warningborder,
    boxrule=1pt,
    left=10pt,
    right=10pt,
    top=8pt,
    bottom=8pt,
    breakable,
    before upper={\faExclamationTriangle\ \textbf{Warning:} }
}

\newtcolorbox{tipbox}{
    colback=tipbg,
    colframe=tipborder,
    boxrule=1pt,
    left=10pt,
    right=10pt,
    top=8pt,
    bottom=8pt,
    breakable,
    before upper={\faLightbulb\ \textbf{Tip:} }
}

\newtcolorbox{notebox}{
    colback=notebg,
    colframe=noteborder,
    boxrule=1pt,
    left=10pt,
    right=10pt,
    top=8pt,
    bottom=8pt,
    breakable,
    before upper={\faInfoCircle\ \textbf{Note:} }
}

% ===== HEADER/FOOTER =====
\pagestyle{fancy}
\fancyhf{}
\fancyhead[L]{\leftmark}
\fancyhead[R]{TwinCAT 3 Lecture Notes}
\fancyfoot[C]{\thepage}
\renewcommand{\headrulewidth}{0.4pt}
\renewcommand{\footrulewidth}{0.4pt}

% ===== HYPERREF SETUP =====
\hypersetup{
    colorlinks=true,
    linkcolor=codeblue,
    filecolor=magenta,
    urlcolor=cyan,
    pdftitle={TwinCAT 3 - ADS Communication},
    pdfauthor={},
    bookmarks=true,
}

% ===== TITLE FORMATTING =====
\titleformat{\section}
    {\normalfont\Large\bfseries}{\thesection}{1em}{}[{\titlerule[0.8pt]}]
\titleformat{\subsection}
    {\normalfont\large\bfseries}{\thesubsection}{1em}{}
\titleformat{\subsubsection}
    {\normalfont\normalsize\bfseries}{\thesubsubsection}{1em}{}

% ===== DOCUMENT INFO =====
\title{
    \vspace{-1cm}
    \textbf{Lecture Notes: TwinCAT 3}\\
    \Large ADS (Automation Device Specification)
}
\author{}
\date{\today}

% ===== BEGIN DOCUMENT =====
\begin{document}

\maketitle

\tableofcontents
\newpage

% ============================================================
\section{Introduction to ADS}
% ============================================================

\textbf{Automation Device Specification (ADS)} is Beckhoff's standardized interface for communication between various software modules within the TwinCAT environment. It utilizes a \textbf{client-server architecture}, where a connection is always initiated by a client to a server.

\subsection{Internal Uses}

ADS facilitates communication between the development environment (XAE) and the runtime (XAR) for reading/writing variables, and it is used internally by the TwinCAT scheduler to coordinate different modules.

\subsection{External Uses}

It allows external devices, such as a Raspberry Pi, Linux machines, or other TwinCAT runtimes, to exchange data with a target PLC.

\subsection{Protocol Details}

ADS can be \textbf{cyclic} (polling) or \textbf{event-based} (subscriptions). It typically uses \textbf{TCP ports 8016 and 48898} for communication and \textbf{UDP port 48899} for broadcasting device searches.

\begin{table}[h]
\centering
\begin{tabular}{@{}lll@{}}
\toprule
\textbf{Port} & \textbf{Protocol} & \textbf{Purpose} \\
\midrule
48898 & TCP & ADS communication \\
8016 & TCP & Secure ADS (TLS) \\
48899 & UDP & Device discovery broadcast \\
\bottomrule
\end{tabular}
\caption{ADS Network Ports}
\end{table}

\begin{figure}[h]
\centering
\begin{tabular}{|c|c|c|}
\hline
\textbf{Client} & $\xrightarrow{\text{ADS Request}}$ & \textbf{Server} \\
(HMI, C\# App, Python) & $\xleftarrow{\text{ADS Response}}$ & (TwinCAT PLC) \\
\hline
\end{tabular}
\caption{ADS Client-Server Architecture}
\end{figure}

% ============================================================
\section{Addressing: AmsNetId and Ports}
% ============================================================

To communicate via ADS, you must identify both the host device and the specific service running on that device.

\subsection{AmsNetId}

This is the unique address of a computer within the TwinCAT network. It consists of \textbf{six octets} (e.g., \texttt{192.168.4.200.1.1}). While often derived from the system's IP address by default, \textbf{the AmsNetId is not an IP address} and can be changed independently.

\begin{notebox}
The AmsNetId format \texttt{x.x.x.x.1.1} is common, where the first four octets match the IP address and the last two are typically \texttt{1.1}. However, this is just a convention---the AmsNetId can be any six-octet value and does not need to correlate with the IP address.
\end{notebox}

\subsection{ADS Ports}

Within a host, different services are assigned specific ports.

\begin{table}[h]
\centering
\begin{tabular}{@{}lll@{}}
\toprule
\textbf{Port} & \textbf{Service} & \textbf{Description} \\
\midrule
10000 & System Service & Logger, Router \\
300 & Real-time I/O & EtherCAT Master/Slaves \\
301 & NC & Numerical Control \\
350 & NC Interpolation & Path interpolation \\
500 & CNC & Computer Numerical Control \\
801--804 & Reserved & Legacy ports \\
851 & PLC Runtime 1 & First PLC instance \\
852 & PLC Runtime 2 & Second PLC instance \\
853 & PLC Runtime 3 & Third PLC instance \\
854 & PLC Runtime 4 & Fourth PLC instance \\
\bottomrule
\end{tabular}
\caption{Common ADS Port Assignments}
\end{table}

\subsection{ADS Message Router}

This component defines communication channels (routes) between different TwinCAT devices. On Windows, you must manually add a route to a remote PLC to establish communication.

\begin{warningbox}
Before an external client can communicate with a TwinCAT PLC, a \textbf{route must be established on both sides}. The Windows TwinCAT system uses the AMS Router GUI, while Linux systems require manual configuration or reverse routing from the PLC side.
\end{warningbox}

% ============================================================
\section{ADS Packet Structure}
% ============================================================

An AMS packet consists of several layers:

\subsection{Packet Layers}

\begin{itemize}[leftmargin=*]
    \item \textbf{AMS TCP Header:} Contains the length of the subsequent data.
    
    \item \textbf{AMS Header:} Contains the \textbf{Target/Source AmsNetId}, \textbf{Target/Source Port}, State Flags (request vs. response), and the \textbf{Command ID} (e.g., Read, Write).
    
    \item \textbf{ADS Data:} The payload, which varies by command. For a ``Read'' command, this includes the \textbf{Index Group}, \textbf{Index Offset}, and length of data.
\end{itemize}

\begin{table}[h]
\centering
\begin{tabular}{@{}lll@{}}
\toprule
\textbf{Command ID} & \textbf{Name} & \textbf{Description} \\
\midrule
0x0001 & Read Device Info & Get device name and version \\
0x0002 & Read & Read data from server \\
0x0003 & Write & Write data to server \\
0x0004 & Read State & Get ADS/device state \\
0x0005 & Write Control & Change ADS state \\
0x0006 & Add Notification & Subscribe to variable changes \\
0x0007 & Delete Notification & Unsubscribe from notifications \\
0x0008 & Device Notification & Notification data from server \\
0x0009 & Read Write & Combined read and write \\
\bottomrule
\end{tabular}
\caption{ADS Command IDs}
\end{table}

\subsection{Index Groups for PLC Access}

\begin{table}[h]
\centering
\begin{tabular}{@{}lll@{}}
\toprule
\textbf{Index Group} & \textbf{Hex} & \textbf{Description} \\
\midrule
16384 & 0x4000 & PLC Memory Area (\%M) \\
16416 & 0x4020 & PLC Data Area (variables) \\
16432 & 0x4030 & Retain Data \\
61443 & 0xF003 & Read by symbol name \\
61449 & 0xF009 & Get symbol handle \\
61450 & 0xF00A & Release symbol handle \\
61451 & 0xF00B & Read by handle \\
61452 & 0xF00C & Write by handle \\
\bottomrule
\end{tabular}
\caption{Common Index Groups for PLC Access}
\end{table}

% ============================================================
\section{Practical Implementation: .NET / C\# (Windows)}
% ============================================================

For Windows-based applications, Beckhoff provides official libraries, including support for \textbf{.NET 6}, which offers cross-platform compatibility.

\subsection{Library Integration}

The easiest method is using the \textbf{NuGet package manager} to install \texttt{Beckhoff.TwinCAT.ADS}.

\subsection{Symbolic Access}

While you can read data via Index Groups (e.g., \textbf{4020} for PLC memory), it is recommended to use \textbf{symbolic paths} (e.g., \texttt{MAIN.nCounter}). This ensures the code remains functional even if variable memory addresses change after recompilation.

\subsubsection*{Example: C\# Symbolic Read}

\begin{lstlisting}[style=csharpstyle]
// C# Console Application: Reading a PLC Variable via ADS
// Required NuGet Package: Beckhoff.TwinCAT.Ads

using System;
using TwinCAT.Ads;

namespace AdsReadExample
{
    class Program
    {
        static void Main(string[] args)
        {
            // Create an ADS client instance
            // This is the main object for ADS communication
            using (AdsClient adsClient = new AdsClient())
            {
                try
                {
                    // Connect to the local TwinCAT runtime on port 851
                    // AmsNetId.Local represents the local machine (127.0.0.1.1.1)
                    // Port 851 is the default port for the first PLC runtime
                    adsClient.Connect(AmsNetId.Local, AmsPort.PlcRuntime_851);
                    
                    Console.WriteLine("Connected to TwinCAT PLC on port 851");
                    
                    // Read a variable using its symbolic path
                    // This is the RECOMMENDED approach as it survives recompilation
                    // The symbol path corresponds to: PROGRAM MAIN, variable nIntToRead
                    short valueRead = (short)adsClient.ReadValue(
                        "MAIN.nIntToRead",  // Symbolic variable path
                        typeof(short)        // Expected data type (INT = short in C#)
                    );
                    
                    Console.WriteLine($"Value of MAIN.nIntToRead: {valueRead}");
                    
                    // Alternative: Read using generic method for type safety
                    // This approach avoids casting and is more robust
                    int nCounter = adsClient.ReadValue<int>("MAIN.nCounter");
                    Console.WriteLine($"Value of MAIN.nCounter: {nCounter}");
                    
                    // Read a REAL (float) value
                    float fTemperature = adsClient.ReadValue<float>("MAIN.fTemperature");
                    Console.WriteLine($"Temperature: {fTemperature:F2} C");
                    
                    // Read a BOOL value
                    bool bMotorRunning = adsClient.ReadValue<bool>("MAIN.bMotorRunning");
                    Console.WriteLine($"Motor Running: {bMotorRunning}");
                    
                    // Read a STRING value (note: TwinCAT strings have fixed length)
                    string sMessage = adsClient.ReadValue<string>("MAIN.sStatusMessage");
                    Console.WriteLine($"Status: {sMessage}");
                }
                catch (AdsException ex)
                {
                    // Handle ADS-specific errors
                    Console.WriteLine($"ADS Error: {ex.Message}");
                    Console.WriteLine($"Error Code: {ex.ErrorCode}");
                }
                catch (Exception ex)
                {
                    // Handle general errors
                    Console.WriteLine($"Error: {ex.Message}");
                }
                finally
                {
                    // Disconnect is handled by 'using' statement
                    Console.WriteLine("Disconnected from PLC");
                }
            }
            
            Console.WriteLine("Press any key to exit...");
            Console.ReadKey();
        }
    }
}
\end{lstlisting}

\begin{lstlisting}[style=csharpstyle]
// Extended Example: Reading from a Remote PLC

using System;
using TwinCAT.Ads;

namespace AdsRemoteReadExample
{
    class Program
    {
        static void Main(string[] args)
        {
            using (AdsClient adsClient = new AdsClient())
            {
                // For remote PLC, specify the AmsNetId explicitly
                // This requires a route to be configured in TwinCAT Router
                string remoteAmsNetId = "192.168.1.100.1.1";
                
                // Parse the AmsNetId string to an AmsNetId object
                AmsNetId targetNetId = AmsNetId.Parse(remoteAmsNetId);
                
                // Connect to remote PLC runtime 1 (port 851)
                adsClient.Connect(targetNetId, AmsPort.PlcRuntime_851);
                
                Console.WriteLine($"Connected to remote PLC: {remoteAmsNetId}");
                
                // Read values from the remote PLC
                int nValue = adsClient.ReadValue<int>("MAIN.nIntToRead");
                Console.WriteLine($"Remote value: {nValue}");
                
                // Write a value to the remote PLC
                adsClient.WriteValue("MAIN.nIntToWrite", 12345);
                Console.WriteLine("Wrote value 12345 to MAIN.nIntToWrite");
            }
        }
    }
}
\end{lstlisting}

\begin{tipbox}
Always use symbolic access (\texttt{ReadValue<T>("MAIN.nVariable")}) instead of raw Index Group/Offset access. Symbolic access automatically resolves variable addresses at runtime, so your application continues to work even after the PLC program is recompiled and variable addresses change.
\end{tipbox}

% ============================================================
\section{Practical Implementation: C++ (Linux)}
% ============================================================

Beckhoff provides an \textbf{open-source ADS client library} on GitHub for non-Windows environments.

\subsection{Compilation}

On Linux (e.g., Debian), the library is compiled from source using tools like \textbf{Meson and Ninja}.

\begin{verbatim}
# Clone the repository
git clone https://github.com/Beckhoff/ADS.git
cd ADS

# Build using Meson and Ninja
meson setup build
ninja -C build

# Install (optional)
sudo ninja -C build install
\end{verbatim}

\subsection{Reverse Routing}

Since Linux lacks the standard TwinCAT AMS Router GUI, you must \textbf{authorize the connection on the PLC side}. This involves adding a route on the target PLC pointing back to the Linux client's IP address and its generated AmsNetId.

\begin{notebox}
On Linux, the ADS library generates an AmsNetId based on the system's hostname. You can find this ID in the library output or specify a custom one. The Windows TwinCAT PLC must have a route configured to this AmsNetId for communication to work.
\end{notebox}

\subsection{C++ API}

The C++ library allows for high-performance reading and writing. It supports overloading assignment operators, making the code very concise.

\subsubsection*{Example: C++ Write Operation}

\begin{lstlisting}[style=cppstyle]
// C++ ADS Client Example: Writing to a TwinCAT PLC Variable
// Requires: Beckhoff ADS open-source library from GitHub
// https://github.com/Beckhoff/ADS

#include <iostream>
#include <string>
#include "AdsLib.h"
#include "AdsVariable.h"

int main()
{
    // Define the target PLC's AmsNetId (six octets)
    // This must match the AmsNetId configured on the TwinCAT system
    static const AmsNetId remoteNetId { 192, 168, 1, 100, 1, 1 };
    
    // Define the target PLC's IP address
    // This is the actual network IP, separate from the AmsNetId
    static const char remoteIpAddress[] = "192.168.1.100";
    
    try
    {
        // Create an ADS device connection
        // Parameters: IP Address, AmsNetId, ADS Port (851 for PLC Runtime 1)
        AdsDevice adsDevice {
            remoteIpAddress,
            remoteNetId,
            AMSPORT_R0_PLC_TC3  // Port 851
        };
        
        std::cout << "Connected to PLC at " << remoteIpAddress << std::endl;
        
        // Create an ADS variable handle for writing
        // Template parameter specifies the data type (int32_t = DINT)
        // The string parameter is the symbolic path in the PLC
        AdsVariable<int32_t> nIntToWrite {
            adsDevice,
            "MAIN.nIntToWrite"
        };
        
        // Write a value to the PLC variable
        // The library overloads the assignment operator for clean syntax
        // This single line performs the entire ADS write operation
        nIntToWrite = 42;
        
        std::cout << "Successfully wrote value 42 to MAIN.nIntToWrite" << std::endl;
        
        // Create another variable for reading back
        AdsVariable<int32_t> nIntToRead {
            adsDevice,
            "MAIN.nIntToRead"
        };
        
        // Read the value - implicit conversion via overloaded operator
        int32_t readValue = nIntToRead;
        std::cout << "Read value from MAIN.nIntToRead: " << readValue << std::endl;
        
        // Example with different data types
        
        // REAL (float)
        AdsVariable<float> fTemperature { adsDevice, "MAIN.fTemperature" };
        float tempValue = fTemperature;
        std::cout << "Temperature: " << tempValue << " C" << std::endl;
        
        // BOOL
        AdsVariable<bool> bMotorRunning { adsDevice, "MAIN.bMotorRunning" };
        bMotorRunning = true;  // Start the motor
        std::cout << "Motor started" << std::endl;
        
        // Writing multiple values in sequence
        AdsVariable<int32_t> nSetpoint { adsDevice, "MAIN.nSetpoint" };
        for (int i = 0; i <= 100; i += 10)
        {
            nSetpoint = i;
            std::cout << "Setpoint: " << i << "%" << std::endl;
            // In real application, add delay here
        }
        
    }
    catch (const AdsException& ex)
    {
        // Handle ADS-specific errors
        std::cerr << "ADS Error: " << ex.what() << std::endl;
        std::cerr << "Error Code: " << ex.errorCode << std::endl;
        return 1;
    }
    catch (const std::exception& ex)
    {
        // Handle general errors
        std::cerr << "Error: " << ex.what() << std::endl;
        return 1;
    }
    
    std::cout << "Program completed successfully" << std::endl;
    return 0;
}
\end{lstlisting}

\begin{lstlisting}[style=cppstyle]
// Alternative C++ Example: Using Low-Level API
// For cases where you need more control over the ADS operations

#include <iostream>
#include "AdsLib.h"

int main()
{
    // Add a route to the remote PLC
    // This is necessary on Linux since there's no AMS Router GUI
    long nPort = AdsPortOpenEx();
    
    if (nPort == 0)
    {
        std::cerr << "Failed to open ADS port" << std::endl;
        return 1;
    }
    
    // Define remote address
    AmsAddr remoteAddr;
    remoteAddr.netId = { 192, 168, 1, 100, 1, 1 };
    remoteAddr.port = AMSPORT_R0_PLC_TC3;  // 851
    
    // Add route (IP address of remote PLC)
    AdsAddRoute(remoteAddr.netId, "192.168.1.100");
    
    // Write a DINT value using Index Group and Offset
    // Index Group 0x4020 = PLC data area
    // Index Offset = memory offset of variable (from symbol info)
    int32_t valueToWrite = 42;
    long nResult = AdsSyncWriteReqEx(
        nPort,
        &remoteAddr,
        0x4020,           // Index Group: PLC Memory
        0x0000,           // Index Offset: variable offset
        sizeof(valueToWrite),
        &valueToWrite
    );
    
    if (nResult != ADSERR_NOERR)
    {
        std::cerr << "Write failed with error: " << nResult << std::endl;
    }
    else
    {
        std::cout << "Write successful" << std::endl;
    }
    
    // Clean up
    AdsPortCloseEx(nPort);
    
    return 0;
}
\end{lstlisting}

\begin{warningbox}
When using the C++ ADS library on Linux, ensure that:
\begin{enumerate}
    \item A route is added on the Windows TwinCAT PLC pointing to your Linux machine's IP and AmsNetId
    \item The firewall allows TCP port 48898 and UDP port 48899
    \item The AmsNetId used in your code matches what's configured in the route
\end{enumerate}
\end{warningbox}

% ============================================================
\section{Summary of Benefits}
% ============================================================

\begin{itemize}[leftmargin=*]
    \item \textbf{Open System:} ADS allows TwinCAT to communicate with almost any language (C\#, C++, Java, Python, Node.js) and platform.
    
    \item \textbf{No Licensing Fees:} Using ADS client libraries is \textbf{free} and does not require additional TwinCAT licenses.
    
    \item \textbf{Full Access:} ADS provides access not just to PLC variables, but also to motion control, I/O, vision, and system analytics.
\end{itemize}

\begin{table}[h]
\centering
\begin{tabular}{@{}llll@{}}
\toprule
\textbf{Platform} & \textbf{Library} & \textbf{Source} & \textbf{License} \\
\midrule
Windows/.NET & Beckhoff.TwinCAT.Ads & NuGet & Free \\
Linux/C++ & ADS Library & GitHub & MIT \\
Python & pyads & PyPI & MIT \\
Node.js & ads-client & npm & MIT \\
Java & AdsToJava & GitHub & Free \\
\bottomrule
\end{tabular}
\caption{ADS Client Libraries by Platform}
\end{table}

\begin{tipbox}
ADS provides a unified interface to access all TwinCAT subsystems. Beyond PLC variables, you can use ADS to:
\begin{itemize}
    \item Read/write NC axis parameters (Port 500/501)
    \item Access I/O terminals directly (Port 300)
    \item Query system diagnostics and logs (Port 10000)
    \item Interface with TwinCAT Vision (Port 2400)
\end{itemize}
\end{tipbox}

% ============================================================
\section*{Quick Reference Card}
\addcontentsline{toc}{section}{Quick Reference Card}
% ============================================================

\subsection*{ADS Addressing Format}

\begin{verbatim}
AmsNetId: x.x.x.x.x.x (six octets, e.g., 192.168.1.100.1.1)
ADS Port: 851 (PLC Runtime 1), 852 (Runtime 2), etc.
Full Address: 192.168.1.100.1.1:851
\end{verbatim}

\subsection*{Common ADS Ports}

\begin{table}[h]
\centering
\begin{tabular}{@{}ll@{}}
\toprule
\textbf{Port} & \textbf{Service} \\
\midrule
300 & I/O (EtherCAT) \\
851--854 & PLC Runtimes 1--4 \\
500 & NC/CNC \\
10000 & System Service \\
\bottomrule
\end{tabular}
\end{table}

\subsection*{C\# Quick Reference}

\begin{lstlisting}[style=csharpstyle]
// Connect
using AdsClient client = new AdsClient();
client.Connect(AmsNetId.Local, AmsPort.PlcRuntime_851);

// Read
int value = client.ReadValue<int>("MAIN.nVariable");

// Write
client.WriteValue("MAIN.nVariable", 42);
\end{lstlisting}

\subsection*{C++ Quick Reference}

\begin{lstlisting}[style=cppstyle]
// Connect
AdsDevice device { "192.168.1.100", {192,168,1,100,1,1}, 851 };

// Create variable handle
AdsVariable<int32_t> myVar { device, "MAIN.nVariable" };

// Read
int32_t value = myVar;

// Write
myVar = 42;
\end{lstlisting}

\vfill

\begin{center}
\textit{Last Updated: \today}
\end{center}

\end{document}
